\chapter{Proportionnalité, ratios et pourcentages}
\label{corriges:07}
\section*{Exercice 7.1 -- Reconnaître le modèle}
Dans le premier, $y/x=2{,}5$ à chaque fois : la formule $y=2{,}5x$ est proportionnelle sur ces données. Dans le second, les quotients diffèrent ; une formule compatible est $y=2x+3$, affine et non proportionnelle. Un nombre fini de données ne détermine cependant pas un modèle unique pour toutes les autres entrées.

\section*{Exercice 7.2 -- Partager et mélanger}
Le partage comporte $2+3+4=9$ parts de $40$ euros : $80$, $120$ et $160$ euros.
Le mélange comporte $6$ parts ; le sirop représente $1/6$ du total, donc $1{,}8/6=0{,}3\unite L=300\unite{mL}$.

\section*{Exercice 7.3 -- Évolutions successives}
Le prix devient $80\times1{,}15\times0{,}8=73{,}60$ euros. Le coefficient global est $0{,}92$, soit une baisse de $8\,\%$. Pour revenir à $80$, le coefficient est $80/73{,}6=25/23$, donc le taux vaut $2/23\approx0{,}08696$, soit environ $8{,}70\,\%$. Une hausse de $8\,\%$ ne suffirait pas.

\section*{Exercice 7.4 -- Échelle et précision}
$7{,}2\times50000=360000\unite{cm}=3{,}6\unite{km}$. Un millimètre sur la carte correspond à $50000\unite{mm}=50\unite m$ en réalité. Cela donne un ordre de grandeur de l'incertitude de lecture, sans prétendre inclure toutes les erreurs de la carte.
